\begin{textsample}{POS Dim 3 – mt – Score 30.00 – mt/mt\_inf\_17.txt}  \label{ex:f3_pos_073}
2.2.4 As Possibilidades de Experiências com \textbf{Bebês}, \textbf{Crianças} Bem Pequenas e \textbf{Crianças} Pequenas Os \textbf{bebês} aprendem a mobilizar as pessoas e a comunicar -se por meio de \textbf{gestos} e expressões, sobretudo na interação com os responsáveis pelos \textbf{cuidados} com sua \textbf{higiene} e alimentação.
Estas \textbf{crianças} ao experienciarem ações com seu corpo, \textbf{gestos} e movimentos, manifestam sensações, se \textbf{deparam} com os desafios corporais como engatinhar, arrastar, ficar de \textbf{pé}, caminhar, subir, descer, correr, rolar, pular, mexer, encaixar e tocar.
A ampliação de tais experiências ocorre conforme elas imitam as expressões, movimentos e falas na interação com seus \textbf{parceiros} \textbf{adultos} ou \textbf{infantis}.
As \textbf{crianças} bem \textbf{pequenas} possuem maiores desenvolvimento motor, afetivo, cognitivo, linguístico e social.
A exploração de objetos pode continuar a acontecer de forma livre para a \textbf{criança}, o faz-de-conta precisa ser potencializado de modo que a \textbf{criança} ao se relacionar com os diferentes, integra -se e expressa as ações dos diversos \textbf{personagens} por meio dos \textbf{gestos} e movimentos.
As \textbf{crianças} \textbf{pequenas}, na progressão do desenvolvimento das expressões corporais e movimentos enriquecem suas possibilidades à medida que exploram o ambiente, principalmente quando fazem uso de referenciais externos para orientar sua experiência corporal em determinado espaço.
Esse campo de experiência, conforme Farias e Salles ( 2012 ) é uma necessidade das \textbf{crianças} no seu processo de inserção e produção cultural, tornando fundamental a exploração de variadas possibilidades, auxiliando à constituição dos sujeitos.
Nesse trabalho, as \textbf{crianças} tomam consciência de seus limites e possibilidades, tanto para proteger sua integridade física, quanto para ter \textbf{cuidado} com o corpo do outro.
O contínuo trabalho com este campo de experiências nos diferentes grupos \textbf{etários}, é fundamental no planejamento curricular cotidiano, objetivando o desenvolvimento integral da \textbf{criança}.
Assim, se faz necessário que dentre as diversas possiblidades, vivenciem experiências de : - Apreciar o teatro de bonecos e fantoches, teatro feito com sombras e as manifestações teatrais com animação de objetos ; - Movimentar -se ritmicamente ao \textbf{som} de músicas de diferentes gêneros ; - Perceber os \textbf{sons} do ambiente e reagir a \textbf{sons} e músicas ; - Reconhecer suas músicas preferidas, acompanhando -as por meio de movimento corporal ; - Produzir \textbf{sons} \textbf{batendo}, sacudindo, chacoalhando, objetos sonoros e instrumentos musicais diversos, usando o próprio corpo e a voz ; - Explorar as qualidades sonoras ( intensidade, duração, timbre, altura ) de objetos e instrumentos musicais diversos mesmos sem reconhecê -las convencionalmente ; - Dançar, criando movimentos ; - Explorar a presença do silêncio como valorização do \textbf{som} ; - Participar de \textbf{brincadeiras} musicais, de \textbf{roda} ou de danças circulares ; - Apreciar músicas instrumentais e diferentes expressões a cultura musical brasileira, bem como outras culturas ; - Possuir autonomia para orientar -se corporalmente com relação a frente, atrás, no alto, em cima, etc ; - Construir, com o auxílio do professor, \textbf{brinquedos} com sucata ; - Ampliar a imitação de \textbf{gestos}, postura e vocalizações de modelos ( \textbf{adultos}, \textbf{crianças}, animais ou \textbf{personagens} de histórias ) ; - Ter presteza e autonomia na manipulação e exploração de diferentes objetos ; - Discriminar e nomear partes do próprio corpo e do outro ; - Controlar gradualmente o próprio movimento, ajustando suas habilidades, as diferentes situações das quais participa ; - Nomear as características e funções das diferentes partes do corpo ; - Apreciar apresentações de danças de diferentes gêneros e outras expressões da cultura corporal ; - Explorar as possibilidades de se expressar, comunicar, interagir intencionalmente com diferentes \textbf{parceiros} pelo movimento ; - Criar e reproduzir coreografias individualmente e em grupo ; - Criar \textbf{brincadeiras} corporais a partir de repertório aprendido ; - Resolver problemas ocorridos em um jogo discutindo regras ; - \textbf{Brincar} de \textbf{esconder} com cabaninhas, lençóis, labirintos ( riscos no chão com barbantes, giz ), \textbf{texturas}, \textbf{areia}, dentre outros ; - Engatinhar, empurrar, rolar, andar, pular, correr, saltar ; - \textbf{Dramatizar} ; - Utilizar os movimentos da \textbf{mão} para rasgar, amassar, apertar, pinçar, empurrar e cortar com tesoura ; - \textbf{Manusear} objetos diversos ( \textbf{lápis}, pincel, giz de cera, tesoura ) ; - Realizar movimentos de preensão, encaixe e lançamento ; - Lançar objetos no espaço a uma determinada distância, coordenando a força necessária para realizar o movimento ; - Ser incentivada e estimulada para executar as ações de sentar sozinha, ficar de \textbf{pé} e andar ; - Apanhar objetos colocados a determinada altura ; - Experimentar odores, sabores, \textbf{texturas} ; - Explorar os cinco sentidos, \textbf{gestos} e emoções ; - Ter \textbf{cuidado} com o seu corpo – higienização, alimentação \textbf{conforto} e aparência ; - \textbf{Brincar} nos espaços externos e internos da \textbf{instituição}, com ou sem obstáculos, desafiando o uso dos diferentes \textbf{gestos} e movimentos corporais ; - Ter experiências de situações que lhes exijam a construção da sua própria identidade e auto-imagem ; - Explorar as habilidades motoras básicas como rolar, andar, correr, pular, dançar, rasgar, recortar ; - Vivenciar diferentes sensações, \textbf{percepções}, emoções ; - Expressar sentimentos, sensações, \textbf{percepções} visual, olfativa, gustativa, auditiva, tátil e cenestésica, nas situações do cotidiano ; - \textbf{Brincar} os diversos sabores, cores, imagens, cheiros, \textbf{texturas}, consistência, temperaturas e etc. ; - Vivenciar \textbf{brincadeiras} e jogos corporais do repertório cultural como : amarelinha, coelhinho sai da toca, \textbf{roda}, boliche, corda, com bolas, bambolê, pneus, obstáculos e outros ; - \textbf{Dramatizar} mediante a situações de imitação e criação de \textbf{personagens}, cenários, tramas, nas \textbf{brincadeiras} de faz de conta.
Assim, a partir das diversas vivências com o corpo, \textbf{gestos} e movimentos, a \textbf{criança} estabelece relações com as pessoas, com os objetos, com os espaços e o tempo, reconhecem seus limites, as possibilidades do seu corpo permitindo de maneira gradativa a apropriação da sua consciência corporal, constituindo sua subjetividade e a ampliação do seu ser e estar no mundo de forma plena e significativa, como fonte de prazer, cultura e aprendizagem.

% matched lemmas: adulto, areia, bater, bebê, brincadeira, brincar, brinquedo, conforto, criança, cuidado, cuidar, deparar, dramatizar, esconder, etário, gesto, higiene, infantil, instituição, lápis, manusear, mão, parceiro, pequeno, percepção, personagem, pé, roda, som, textura
\end{textsample}
